\section{Summary and Conclusion}
This thesis presented the comprehensive design, comparative analysis, and experimental validation of a flatwise-wound planar transformer engineered to resolve the fundamental thermal bottlenecks in high-power-density Solid State Transformers (SSTs). Traditional planar transformers predominantly utilize an edgewise winding architecture, which inherently traps heat due to the high series thermal resistance of the inter-turn dielectric insulation. To overcome this critical limitation, this study proposed a structural shift to a flatwise orientation, establishing a highly conductive parallel thermal path that grants every individual copper turn direct, orthogonal access to the external cooling boundaries.

The theoretical advantages of this geometric reconfiguration were first evaluated through comprehensive 3D Multi-Physics Finite Element Analysis (FEA). Baseline simulations confirmed that the flatwise orientation inherently provides a superior thermal path by significantly lowering the peak operating temperature and eliminating localized hotspots under both natural and forced-air cooling conditions compared to the conventional edgewise equivalent. 

To definitively validate these findings, hardware prototypes of both winding configurations were fabricated and subjected to continuous, full-load experimental evaluation in a 1~kW, 28.5~kHz bidirectional CLLC resonant converter. The hardware testing conclusively validated the core hypothesis of this research. Under continuous steady-state operation, the conventional edgewise design suffered from severe internal-to-external thermal gradients, resulting in trapped heat and the rapid onset of thermal runaway. In stark contrast, the flatwise design consistently maintained safe operation, achieving significantly lower thermal gradients and reduced peak winding temperatures across all tested ambient conditions. By narrowing the temperature gradients between the internal copper windings and the core exterior, the experimental data proves that the flatwise architecture successfully establishes a fundamentally lower thermal resistance path. This enables highly efficient heat extraction directly from the core window, ultimately breaking the restrictive dielectric-thermal trade-off.

Finally, the scalability of the flatwise winding architecture was examined under controlled parameters. A geometric scalability analysis evaluated across varying core window aspect ratios at 1~kW revealed that the flatwise structure guarantees exceptional heat extraction even in high or low-profile physical geometries, establishing that the configuration is versatile in terms of spatial geometry. Furthermore, a power scalability analysis utilizing the Area-Product ($AP$) method evaluated the configurations at 10~kW and 100~kW operating limits maintaining a square window area. The results demonstrated that as the power rating and corresponding interwinding distances scale up, the thermal divergence between the two structures amplifies drastically. The flatwise-wound sturcture is able to reduce high-power peak temperatures, delivering exceptional thermal stability in exchange for only a marginal increase in total power loss at almost the same foot print with its conventional counterpart.

Ultimately, the comprehensive findings of this study validate the flatwise-wound planar transformer as a structurally superior candidate for compact, high-reliability power converters. By providing massive thermal headroom with only a marginal trade-off in high-frequency winding loss, this architecture successfully unlocks new limits for volumetric miniaturization in next-generation electrical distribution infrastructure.

\section{Recommendations and Future Work}
Building upon the established results of this study, a few key areas for recommendation and future work are suggested as follows:

\begin{itemize}
    \item \textbf{Model Verification and Optimization:} While the 3D Multi-Physics FEA provided a robust theoretical foundation for the thermal advantages of the flatwise orientation, further refinement of the analytical models is recommended. Developing faster analytical thermal and electromagnetic models will provide critical insights for design optimization, especially for larger, high-power applications where extensive 3D FEA parametric sweeps become computationally prohibitive.
    
    \item \textbf{Parasitics Modeling and Mitigation:} The proposed transformer utilizes rectangular windings. This geometry implies that parasitic capacitance will be a dominant factor, as rectangular windings feature large overlapping surface areas. A direct analogy can be drawn to PCB-based planar transformers, where large overlapping areas between adjacent copper layers inherently generate significant inter-turn and inter-winding capacitances. Future work should prioritize the development of an accurate high-frequency analytical model to extract these distributed parasitic elements. Understanding and quantifying these capacitive effects is strictly necessary to accurately predict common-mode electromagnetic interference (EMI) and mitigate unintended high-frequency ringing, which are critical steps before integrating the transformer into high-voltage, high-frequency SST applications as an example.
    
    \item \textbf{Further Scalability Analysis and High-Voltage Experimentation:} As this study successfully validated the flatwise architecture on a scaled-down 1~kW prototype, future work must formally evaluate the structural and thermal viability of this design toward the ultimate 100~kW, 1000~V target hardware. Transitioning to the 1~kV operating regime necessitates strict adherence to international safety and insulation standards, such as IEC 60664 and IEC 62477. As established in Chapter 6, future physical iterations must incorporate comprehensive insulation coordination, explicitly meeting rigorous creepage and clearance distance requirements. Experimental validation at this scale will be crucial to ensure that the physical structural modifications required for high-voltage compliance do not inadvertently compromise the established thermal advantages of the flatwise winding configuration.
\end{itemize}